\subsection{测度的提升}

命题 8. --- 设 T 和 X 是两个拓扑空间，\(\pi\) 是从 T 到 X 的映射。

\begin{enumerate}
\def\labelenumi{\alph{enumi})}
\item
  设 \(\nu\) 是 X 上的有界测度。要使存在 T 上的测度 \(\mu\)，使得 \(\pi\) 是 \(\mu\) -逆紧的且 \(\pi(\mu) = \nu\)，必要且充分的条件是：对于每个数 \(\varepsilon > 0\)，存在紧集 \(\mathrm{K}_{\varepsilon} \subset \mathrm{T}\)，使得 \(\pi\) 在 \(\mathrm{K}_{\varepsilon}\) 上的限制是连续的，并且 \(\nu^{\bullet}(\mathrm{X} - \pi(\mathrm{K}_{\varepsilon})) < \varepsilon\)。
\item
  假设 \(\pi\) 是单射；令 \(\mu\) 和 \(\mu'\) 是 T 上的两个测度，使得 \(\pi\) 对 \(\mu\) 和 \(\mu'\) 都是逆紧的，并且 \(\pi(\mu) = \pi(\mu')\)。则 \(\mu = \mu'\)。
\end{enumerate}

所述条件 \(a\) 是必要的。事实上，若 \(\pi\) 是 \(\mu\) -逆紧的且 \(\pi(\mu) = \nu\)，则关系 \(\mu^{\bullet}(1) = \nu^{\bullet}(1) < +\infty\) 蕴含 \(\mu\) 有界。将 §1, No.~2 的 Prop. 2 应用于函数 1，得到 T 中存在紧子集 K，使得 \(\mu^{\bullet}(\mathrm{T} - \mathrm{K}) < \varepsilon/2\)。由于 \(\pi\) 是 \(\mu\) -可测的，存在紧集 \(\mathrm{K}_{\varepsilon} \subset \mathrm{K}\)，使得 \(\pi\) 在 \(\mathrm{K}_{\varepsilon}\) 上的限制连续，并且 \(\mu^{\bullet}(\mathrm{K} - \mathrm{K}_{\varepsilon}) < \varepsilon/2\)。于是（No.~3, Prop. 4）

\[
\nu^ {\bullet} \big (\mathrm{X} - \pi (\mathrm{K} _ {\varepsilon}) \big) = \mu^ {\bullet} \big (\mathrm{T} - \pi^{-1} \big (\pi (\mathrm{K} _ {\varepsilon}) \big) \big) <   \varepsilon .
\]

为证明条件是充分的，我们先处理一个特殊情形。

引理 1. --- 设 U 和 V 是两个紧空间，h 是从 U 到 V 的连续满射。于是映射 \(\lambda \mapsto h(\lambda)\) 从 \(\mathcal{M}_{+}(U)\) 到 \(\mathcal{M}_{+}(V)\) 是满射。

事实上，令 \(a\) 为线性映射 \(f \mapsto f \circ h\)，它从 \(\mathcal{C}(\mathrm{V})\) 映到 \(\mathcal{C}(\mathrm{U})\)；由于 \(h\) 是满射，\(a\) 是从 \(\mathcal{C}(\mathrm{V})\) 到子空间 \(\mathrm{H}\)（\(\mathcal{C}(\mathrm{U})\) 的子空间）的等距映射。令 \(\theta\) 是 \(\mathrm{V}\) 上的正测度；于是 \(\theta \circ a^{-1}\) 是 \(\mathrm{H}\) 上的连续线性形式，可延拓为 \(\eta\)（定义在 \(\mathcal{C}(\mathrm{U})\) 上）且具有相同范数。于是 \(\eta\) 是 \(\mathrm{U}\) 上的测度，并且 \(\theta(f) = \eta(f \circ h)\) 对所有 \(f \in \mathcal{C}(\mathrm{V})\) 成立，所以 \(\theta = h(\eta)\)。最后，\(\theta(1) = \| \theta \| = \| \eta \|\)，且 \(\theta(1) = \eta(1)\)，故 \(\eta\) 是正的（Ch. V, §5, No.~5, Prop. 9）。

现在证明条件 \(a\) ）的充分性。该条件蕴含存在 \((\mathbf{K}_n)_{n \geqslant 1}\) 的紧子集序列，定义在 \(T\) 中，使得 \(\pi\) 在每个 \(\mathbf{K}_n\) 上的限制连续，并且对于每个 \(n\) 都有 \(\nu^{\bullet}(\mathrm{X} - \pi(\mathrm{K}_n)) < 1/n\)。可以假设序列 \((\mathbf{K}_n)\) 是递增的。置 \(\mathrm{L}_n = \pi(\mathrm{K}_n)\)，并以 \(\nu_n'\) 表示 \(\varphi_{\mathrm{L}_n - \mathrm{L}_{n-1}} \cdot \nu_{\mathrm{L}_n}\) 上的测度 \(\mathrm{L}_n\)，约定 \(\mathrm{L}_0 = \emptyset\)。

由于限制 \(\pi_{\mathrm{K}_n}\) 连续，存在测度 \(\mu_n'\)，它定义在 \(\mathbf{K}_n\) 上，使得 \(\pi_{\mathrm{K}_n}(\mu_n') = \nu_n'\)（Lemma 1）。令 \(\mu_n\) 为 \(\mu_n'\) 在 \(\mathbf{K}_n\) 到 \(\mathbf{T}\) 的典范注入下的像，并令 \(g\) 是 \(\mathcal{F}_{+}(\mathbf{X})\) 中的元素。依次使用 \(\nu\) 集中于 \(\bigcup_{n} \mathrm{L}_n\) 这一事实、§1, No.~5 的 Prop. 4、§1, No.~2 的 Prop. 2、No.~3 的 Prop. 4 以及最后 No.~3 的 Prop. 7，可得

\[
\begin{array}{r l} \nu^ {\bullet} (g) & = \sum_ {n} \nu^ {\bullet} (\varphi_ {\mathrm{L} _ {n} - \mathrm{L} _ {n - 1}} g) = \sum_ {n} \nu_ {n} ^ {\prime   \bullet} (g _ {\mathrm{L} _ {n}}) = \sum_ {n} \mu_ {n} ^ {\prime   \bullet} (g _ {\mathrm{L} _ {n}} \circ \pi_ {\mathrm{K} _ {n}}) \\ & = \sum_ {n} \mu_ {n} ^ {\prime   \bullet} ((g \circ \pi) _ {\mathrm{K} _ {n}}) = \sum_ {n} \mu_ {n} ^ {\bullet} (g \circ \pi). \end{array}
\]

在这个公式中取 \(g = 1\)，可见族 \((\mu_n)\) 可求和，且其和是一个有界测度 \(\mu\)（§1, No.~7, Prop. 7）。根据 No.~3 的 Prop. 5，对于所有 \(\pi\)，映射 \(\mu_n\) 是 \(n\) -可测的，因为 \(\pi_{\mathrm{K}_n}\) 连续，因而是 \(\mu_n'\) -可测的；由此 \(\pi\) 是 \(\mu\) -可测的（§1, No.~7, Prop. 8），又因 \(\mu\) 有界而是 \(\mu\) -逆紧的。上述关系于是证明测度 \(\pi(\mu)\) 与 \(\nu\) 具有相同的本质上积分，故二者相等（§1, No.~2, Cor. of Prop. 2）。

最后，假设 \(\pi\) 是单射，并证明 \(b\)。令 \(f\) 是 \(\mathcal{F}_{+}(\mathrm{T})\) 中的元素；由于 \(\pi\) 是单射，存在函数 \(g \in \mathcal{F}_{+}(\mathrm{X})\)，使得 \(f = g \circ \pi\)，并且置 \(\nu = \pi(\mu) = \pi(\mu')\) 后，由 No.~3 的 Prop. 4 有

\[
\mu^ {\bullet} (f) = \mu^ {\bullet} (g \circ \pi) = \nu^ {\bullet} (g) = \mu^ {\prime \bullet} (g \circ \pi) = \mu^ {\prime \bullet} (f).
\]

因此测度 \(\mu\) 和 \(\mu'\) 具有相同的本质上积分，这蕴含二者相等（§1, No.~2, Cor. of Prop. 2）。

注记。--- 假设 \(\pi\) 是单射。设 \(\theta\) 是复测度，\(\pi\) 是 \(\theta\) -逆紧的且 \(\pi(\theta)=0\)；则 \(\theta=0\)。事实上，将 \(\theta\) 分解为实部和虚部后，可以归结为 \(\theta\) 为实测度的情形。此时有 \(\pi(\theta^{+})=\pi(\theta^{-})\)，因而 \(\theta^{+}=\theta^{-}\)（Prop. 8），最后 \(\theta=0\)。

下面是 Prop. 8 的条件 \(a)\) 总是满足的一个重要情形。

命题 9. --- 设 T 是苏斯林空间（GT, IX, §6, No.~2, Def. 2），X 是 Hausdorff 空间，\(\pi\) 是从 T 到 X 的连续满射，\(\nu\) 是 X 上的有界测度。则存在 T 上的有界测度 \(\mu\)，使得 \(\pi(\mu) = \nu\)。

这些假设显然蕴含 \(\mathbf{X}\) 是苏斯林空间。

考虑集合函数 \(c: \mathrm{A} \mapsto \nu^{\bullet}(\pi(\mathrm{A}))\)，它定义在 \(\mathfrak{P}(\mathrm{T})\) 上。关系 \(\mathrm{A} \subset \mathrm{B}\) 蕴含 \(c(\mathrm{A}) \leqslant c(\mathrm{B})\)；若 \((\mathrm{A}_n)\) 是 \(\mathrm{T}\) 的子集的递增序列，且 \(\mathrm{A} = \bigcup_{n \in \mathbf{N}} \mathrm{A}_n\)，则 \(c(\mathrm{A}) = \sup_n c(\mathrm{A}_n)\)，这是因为 \(\nu^{\bullet}\) 是负荷。最后，令 \(\mathrm{A} \subset \mathrm{T}\)，令 \(\varepsilon\) 是一个 \(> 0\) 的数；取 \(\mathrm{G}\) 为 \(\mathrm{X}\) 中包含 \(\pi(\mathrm{A})\) 的开子集，使得 \(\nu^{\bullet}(\mathrm{G}) \leqslant \nu^{\bullet}(\pi(\mathrm{A})) + \varepsilon\)（§1, No.~9, Prop. 13）；\(\mathrm{H} = \overline{\pi}^{-1}(\mathrm{G})\) 中的开子集 \(\mathrm{T}\) 包含 \(\mathrm{A}\)，并且 \(c(\mathrm{H}) \leqslant c(\mathrm{A}) + \varepsilon\)。因此函数 \(c\) 是 \(\mathrm{T} (\mathrm{TG},\mathrm{IX},\S 6,\mathrm{No.~10},\mathrm{Def.~9})^{(1)}\) 上的右连续容度，而可容性定理（loc. cit., Th. 6）蕴含等式 \(c(\mathrm{T}) = \sup_{\mathrm{K}}c(\mathrm{K})\)，其中 \(K\) 遍历 \(T\) 的紧子集。于是 Prop. 8 蕴含所需测度 \(\mu\) 的存在性。

